\begin{lemma}[The order-nine and height-five parameters]
Let $G$ be a finite loopless multigraph, let $k$ be a natural number, and
suppose that $G$ is critical, $\chi'(G)=k+1$, $\Delta(G)+1\leq k$, and
every vertex is incident with an edge.  Suppose also that $G$ is not
elementary and $t(G)<11$.  If $(T,e_0,\varphi,N)$ is a balanced normal
globally maximum triple, then
\[
 t(G)=9,\qquad h(G)=5,\qquad T.\mathsf{length}=8,
 \qquad N.\mathsf{height}=5.
\]
\end{lemma}
\begin{proof}
If $h(G)<5$, \noderef{ScheideSmallHeightElementary} would make $G$
elementary.  Hence $5\leq h(G)$.  Similarly, if
$t(G)-4<h(G)$, \noderef{ScheideLargeHeightElementary} would make $G$
elementary, so $h(G)\leq t(G)-4$.

By \noderef{ScheideGraphHeightOdd}, $h(G)$ is odd.  The order assertion
and oddness assertion of \noderef{ScheideGlobalMaximumOrder}, applied to
the globally maximum field of $N$, say that
\[
 t(G)=|V(T)|\quad\hbox{and}\quad |V(T)|\text{ is odd}.
\]
The injectivity field of \noderef{TreeSequence}, together with the
definition \noderef{TreeSequenceVertexSet}, gives
$|V(T)|=T.\mathsf{length}+1$.  Thus $t(G)$ is odd.  The inequalities
$5\leq h(G)\leq t(G)-4$ give $9\leq t(G)$, while $t(G)<11$ and oddness
give $t(G)=9$.  The same inequalities give $5\leq h(G)\leq5$, hence
$h(G)=5$.  The two cardinality equalities now give
$T.\mathsf{length}=8$.  Finally, the defining height equality in
\noderef{BalancedGlobalMaximumTriple} says
$N.\mathsf{height}=h(G)=5$.
\end{proof}
